\documentclass[border=4pt]{standalone}
\usepackage[siunitx]{circuitikz}

\begin{document}
\begin{circuitikz}[european resistors]
  % Piezoelectric Norton equivalent: generated charge-rate current, sensor
  % capacitance, and leakage share the input node.
  \coordinate (TP0) at (0,0);
  \coordinate (TPIN) at (0,3.5);
  \draw (TP0) node[ground] {}
    to[I] (TPIN);
  \draw[->] (-0.80,1.00) -- (-0.80,2.80);
  \node[font=\small, align=right, anchor=east] at (-1.05,1.75)
    {$i_q=\mathrm{d}q_s/\mathrm{d}t$};
  \draw (1.5,0)
    to[C] (1.5,3.5);
  \node[font=\small, anchor=west] at (1.90,2.10) {$C_P$};
  \draw (3.0,0)
    to[R] (3.0,3.5);
  \node[font=\small, anchor=west] at (3.45,1.75) {$R_P$};
  \draw (TPIN) -- (3.0,3.5);
  \draw (TP0) -- (3.0,0);

  % Align the inverting input with the sensor node so the signal and both
  % feedback paths meet at one unambiguous summing junction.
  \node[op amp, scale=1.25, anchor=center] (A1) at (7.1,2.90) {};
  \draw (3.0,3.5) -- (A1.-);
  \draw (A1.+) -- (A1.+ |- 0,0) node[ground] {};

  % Feedback capacitance converts charge to voltage; RF supplies a DC path.
  \coordinate (FBOUT) at (9.7,2.90);
  \draw (A1.out) -- (FBOUT)
    node[ocirc, label={[font=\small]right:{TPO: $v_o$}}] {};
  \draw (A1.-) -- (A1.- |- 0,5.80)
    to[C] (FBOUT |- 0,5.80)
    -- (FBOUT);
  \draw (A1.-) -- (A1.- |- 0,4.55)
    to[R] (FBOUT |- 0,4.55)
    -- (FBOUT);
  \node[font=\small] at (8.0,6.5) {$C_F$};
  \node[font=\small] at (8.0,4.95) {$R_F$};

  \node[font=\small\bfseries] at (1.5,6.15)
    {piezoelectric terminal equivalent};
  \node[circ] at (0,3.5) {};
  \node[circ] at (0,0) {};
  \node[circ] at (3.0,3.5) {};
  \node[circ] at (3.0,0) {};
  \node[circ] at (A1.-) {};
  \node[circ] at (FBOUT) {};
\end{circuitikz}
\end{document}
